
On dispose des données suivantes, issues de l'échantillonnage d'une fonction $f(x)$ :

\[
(x_i, f(x_i)) = 
(2,\ 4016),\quad 
(4,\ 1296),\quad 
(5,\ 635),\quad 
(6,\ 256),\quad 
(7,\ 51).
\]

\begin{enumerate}    
    \item Utiliser les trois premiers points pour déterminer le polynôme d'interpolation de degré 2, puis en déduire une approximation de $f(3)$.
    
    \item Utiliser l'ensemble des cinq points pour construire le polynôme d'interpolation de degré 4, puis en déduire une approximation de $f(3)$.
    
    \item Proposer plusieurs méthodes pour approcher l'intégrale
    \[
    \int_2^7 f(x)\, \mathrm{d}x
    \]
    à l'aide des données disponibles. 
\end{enumerate}